\subsection{Types}
\label{sec:mailbox-types}

\begin{figure}[t]
    \small
  \centering
    \begin{tabular}{lccl}
      Mailbox types (\rocqlink{https://edgardschi.github.io/mailbox-types-rocq/MailboxTypes.Types.html\#MType}) \; & $J, K$ \; & $::=$ \; & $\mbtin{E} \mid \mbtout{E}$\\
      Base types (\rocqlink{https://edgardschi.github.io/mailbox-types-rocq/MailboxTypes.Types.html\#BType}) \; & $C$ \; & $::=$ \; & $\unittype \mid \booltype$\\
      Types (\rocqlink{https://edgardschi.github.io/mailbox-types-rocq/MailboxTypes.Types.html\#TType}) \; & $T, U$ \; & $::=$ \; & $C \mid J$\\
    \end{tabular}
  \caption{The syntax of types.}
  \label{fig:type-syntax}
\end{figure}

A mailbox type (Figure \ref{fig:type-syntax}) consists of a capability and a pattern.
A capability is either \emph{input} $\mbtin{}$ or \emph{output} $\mbtout$.
Together, capability and pattern describe what messages are expected to be received from ($\mbtin$) or stored in ($\mbtout$) the mailbox.
To build an intuition for mailbox types, consider the following examples:
A mailbox with mailbox type $\mbtin{\mbmsg{m}}$ gives a process the right to receive a message \msg{m} from that mailbox,
while a mailbox with type $\mbtout{\mbmsg{m}}$ requires the process to send a message $\msg{m}$ to that mailbox.
% must store a $\msg{m}$ message in its mailbox, while a process with type $\mbtin{\mbmsg{m}}$ expects to receive a $\msg{m}$ message.
For a mailbox with mailbox type $\mbtout{(\mbmsg{m} \mbchoice \mbmsg{n})}$, a process may choose whether to send message $\msg{m}$ or $\msg{n}$,
whereas a mailbox with type $\mbtin{(\mbmsg{m} \mbchoice \mbmsg{n})}$ allows a process to receive either message $\msg{m}$ or $\msg{n}$ from that mailbox.
% A process with mailbox type \mbtout{(\mbmsg{m} \mbchoice \mbmsg{n})} is able to choose whether to store message $\msg{m}$ or $\msg{n}$,
% whereas a process using type \mbtin{(\mbmsg{m} \mbchoice \mbmsg{n})} needs to be able to receive both messages.
The type $\mbtin{(\mbmsg{m} \mbcomp \mbmsg{n})}$ is associated to a mailbox that
allows a process to receive both $\msg{m}$ and $\msg{n}$, while a mailbox with mailbox type $\mbtout{(\mbmsg{m} \mbcomp \mbmsg{n})}$ requires a process to send both $\msg{m}$ and $\msg{n}$ to the mailbox, in whichever order.
Finally, a mailbox with type $\mbtout{(\mbstar{\mbmsg{m}})}$ allows a process
to send an arbitrary number of $\msg{m}$ messages to the mailbox,
while $\mbtin{(\mbstar{\mbmsg{m}})}$ describes a mailbox from which a process is able to receive an arbitrary number of $\msg{m}$ messages.
% is able to decide how many $\msg{m}$ messages its mailbox stores, while $\mbtin{(\mbstar{\mbmsg{m}})}$ describes a mailbox which must be able to receive an arbitary amount of $\msg{m}$ messages.

In addition to mailbox types, we also include \emph{base types}.
While the unit type $\unittype$ is required by the type system, additional base types can be added.
For demonstration purposes we also include a Boolean type $\booltype$.
A type $T$ is then either some base type or a mailbox type.

% To deal with the challenges of name aliasing and cyclic dependencies, described in Section \ref{sec:pat-intro}, Pat's solution is to use quasi-linear types.
% To this end, mailbox types with \emph{usage annotations} are used.
% Base types are not equipped with usage, as they do not suffer from the problems as mailbox types do.
% There are two usage annotations: \emph{second class} ($\usecondclass$) and \emph{returnable} ($\ureturnable$).
% Values of second class types can be used several times within the scope they are defined, however cannot leave that scope.
% On the other hand, values of returnable types are consumed once and are allowed to appear in the return type of an expression.
% Usage types ensure the following properties \cite{fowlerSpecialDeliveryProgramming2023}:
% \begin{enumerate}
%   \item There exists only a single returnable reference for each mailbox name in a process
%   \item It is only possible to rename returnable references, thus avoiding aliasing problems with mailbox names
%   \item A returnable reference is the final use of a mailbox name in a process
% \end{enumerate}

% For modifying usage-annotated types, we introduce several operators acting on such types, as seen in Figure \ref{fig:type-syntax}.

% \begin{linkdef}{https://edgardschi.github.io/mailbox-types-rocq/MailboxTypes.Types.html\#type_def}[Type Operators]
%   The operators $\secondclassop{\;\cdot\;}$ and $\returnableop{\;\cdot\;}$ transform a type into a second class or returnable type, respectively.
%   Additionally, the operator $\eraseop{\;\cdot\;}$ removes usage annotations from a type.
%
%   \vspace{-.4cm}
%   \begin{minipage}[t]{.32\linewidth}
%     \begin{align*}
%       \secondclassop{C} &= C\\
%       \secondclassop{T} &= T^\usecondclass\\
%       \secondclassop{T^\eta} &= T^\usecondclass
%     \end{align*}
%   \end{minipage}
%   \begin{minipage}[t]{.32\linewidth}
%     \begin{align*}
%       \returnableop{C} &= C\\
%       \returnableop{T} &= T^\ureturnable\\
%       \returnableop{T^\eta} &= T^\ureturnable\\
%     \end{align*}
%   \end{minipage}
%   \begin{minipage}[t]{.32\linewidth}
%     \begin{align*}
%       \eraseop{C} &= C\\
%       \eraseop{T^\eta} &= T
%     \end{align*}
%   \end{minipage}
% \end{linkdef}

% \subsubsection{Subtyping}
% \label{sec:subtyping}
%
% \begin{figure}[t]
%   \centering
%   \begin{center}
%         \begin{prooftree}
%             \infer0{C \subtype C}
%         \end{prooftree}
%         \quad
%         \begin{prooftree}
%             \hypo{E \mbpinc F}
%             \infer1{\mbtin{E} \subtype \mbtin{F}}
%         \end{prooftree}
%         \quad
%         \begin{prooftree}
%             \hypo{F \mbpinc E}
%             \infer1{\mbtout{E} \subtype \mbtout{F}}
%         \end{prooftree}
%     \end{center}
%   % \begin{minipage}[t]{0.6\textwidth}\vspace{0pt}
%   %   \textbf{Subtyping} \hfill \fbox{$A \subtype B$}
%   %   \begin{center}
%   %         \begin{prooftree}
%   %             \infer0{C \subtype C}
%   %         \end{prooftree}
%   %         \quad
%   %         \begin{prooftree}
%   %             \hypo{E \mbpinc F}
%   %             \hypo{\eta_1 \subtype \eta_2}
%   %             \infer2{\mbtin{E^{\eta_1}} \subtype \mbtin{F^{\eta_2}}}
%   %         \end{prooftree}
%   %         \quad
%   %         \begin{prooftree}
%   %             \hypo{F \mbpinc E}
%   %             \hypo{\eta_1 \subtype \eta_2}
%   %             \infer2{\mbtout{E^{\eta_1}} \subtype \mbtout{F^{\eta_2}}}
%   %         \end{prooftree}
%   %     \end{center}
%   % \end{minipage}\hfill
%   % \begin{minipage}[t]{0.37\textwidth}\vspace{0cm}
%   %   \textbf{Usage subtyping} \hfill \fbox{$\eta_1 \subtype \eta_2$}
%   %   \begin{align*}
%   %     \eta \subtype \eta\\
%   %     \ureturnable \subtype \usecondclass
%   %   \end{align*}
%   % \end{minipage}
%
%   \caption{The semantics of subtyping}
%   \label{fig:subtyping-definition}
% \end{figure}

% \todo{ES}{Maybe move this entire subsection into the Pat section. We don't really need the definition of subtyping here.}
% We now turn to one of the most complex aspects of mailbox types: \emph{subtyping}.
% It is needed in mailbox typing to transform types into suitable forms, for example for guard expressions.
% Subtyping (Figure \ref{fig:subtyping-definition}) is \emph{covariant} for mailbox types with receive capability and \emph{contravariant} for mailbox types with send capability \cite{deliguoroMailboxTypesUnordered2018}.
% Intuitively, this means that a receiving mailbox can be replaced by another mailbox that can receive more messages, and a sending mailbox can be replaced by a another mailbox that can send fewer messages.
% As an example, a mailbox with type $\mbtin{\mbmsg{m}}$ can safely be replaced with one having type $\mbtin{(\mbmsg{m} \mbchoice \mbmsg{n})}$, since it is allowed to receive more messages.
% Similarly, for a mailbox with type $\mbtout{(\mbmsg{m} \mbchoice \mbmsg{n})}$ can safely be replaced with one of type $\mbtout{\mbmsg{m}}$, since it sends a smaller amount of messages.

% Before defining subtyping on mailbox types, we have to define subtyping for usage-annotations.

% \begin{linkdef}{https://edgardschi.github.io/mailbox-types-rocq/MailboxTypes.Types.html\#UsageSubtype}[Usage Subtyping]
%   \emph{Usage subtyping} (written $\eta_1 \subtype \eta_2$) is defined by the following axioms:
%   $$\eta \subtype \eta \hskip0.2\textwidth \ureturnable \subtype \usecondclass.$$
% \end{linkdef}

% Subtyping on usage annotations only allows reflexivity, or for returnable types to be subtypes of second class types.
% If we would also also subtyping in the other direction, i.e.\ $\usecondclass \subtype \ureturnable$, it would be possible to escape the restrictions of second class types using subtyping.
% With subtyping defined on usage annotations, we can continue with subtyping on types.

% \begin{linkdef}{https://edgardschi.github.io/mailbox-types-rocq/MailboxTypes.Types.html\#Subtype}[Subtyping]
%   \emph{Subtyping} (written $A \subtype B$) is defined by the rules:
%   \begin{center}
%         \begin{prooftree}
%             \infer0{C \subtype C}
%         \end{prooftree}
%         \qquad
%         \begin{prooftree}
%             \hypo{E \mbpinc F}
%             \hypo{\eta_1 \subtype \eta_2}
%             \infer2{\mbtin{E^{\eta_1}} \subtype \mbtin{F^{\eta_2}}}
%         \end{prooftree}
%         \qquad
%         \begin{prooftree}
%             \hypo{F \mbpinc E}
%             \hypo{\eta_1 \subtype \eta_2}
%             \infer2{\mbtout{E^{\eta_1}} \subtype \mbtout{F^{\eta_2}}}
%         \end{prooftree}
%     \end{center}
% \end{linkdef}

% Base types can only be subtypes of themselves.
% As mentioned previously, subtyping receive capabilities is covariant and subtyping send capabilities is contravariant.
% Thus, in the subtyping rules for $\mbtin$ and $\mbtout$, the pattern inclusion requirement is flipped.

\paragraph{Type Combination.}
\label{sec:type-combiation}
To model the dynamic nature of mailboxes receiving and sending messages, mailbox types should be able to be combined to create a new mailbox type.
Mailbox types need to ensure that sends and receives are matched, i.e.\ a send is required to have a corresponding receive.
For example, given mailbox types $\mbtout{\mbmsg{m}}$ and $\mbtin{(\mbmsg{m} \mbcomp \mbmsg{n})}$,
the resulting type from the type combination should be $\mbtin{\mbmsg{n}}$.

\begin{figure}[t]
    {\small
  \begin{minipage}[t]{0.55\textwidth}
  \textbf{Pattern permutation} (\rocqlink{https://edgardschi.github.io/mailbox-types-rocq/MailboxTypes.Types.html\#PatternEq}) \hfill \fbox{$E \sim F$}
  \begin{center}
        \begin{prooftree}
            \infer0{(E \mbcomp F) \mbpeqca (F \mbcomp E)}
        \end{prooftree}
        \quad
        \begin{prooftree}
          \hypo{F \mbpeqca F'}
          \infer1{E \mbcomp F \mbpeqca E \mbcomp F'}
        \end{prooftree}
    \begin{center}
        \begin{prooftree}
            \infer0{E \mbpeqca E}
        \end{prooftree}
        \quad
        \begin{prooftree}
          \hypo{F \mbpeqca E}
          \infer1{E \mbpeqca F}
        \end{prooftree}
        \quad
        \begin{prooftree}
          \hypo{E \mbpeqca F'}
          \hypo{F' \mbpeqca F}
          \infer2{E \mbpeqca F}
        \end{prooftree}
    \end{center}
  \end{center}
  \begin{center}
        \begin{prooftree}
            \infer0{E \mbcomp (F \mbcomp F') \mbpeqca (E \mbcomp F) \mbcomp F'}
        \end{prooftree}
  \end{center}
  \end{minipage}
  \hfill
  \begin{minipage}[t]{0.44\textwidth}
    \textbf{Type permutation} (\rocqlink{https://edgardschi.github.io/mailbox-types-rocq/MailboxTypes.Types.html\#TypeEq}) \hfill \fbox{$T \sim U$}
  \begin{center}
        \begin{prooftree}
            \infer0{C \mbpeqca C}
        \end{prooftree}
        \qquad
        \begin{prooftree}
          \hypo{E \mbpeqca F}
          \infer1{\mbtout{E} \mbpeqca \mbtout{F}}
        \end{prooftree}
  \end{center}
  \begin{center}
        \begin{prooftree}
          \hypo{E \mbpeqca F}
          \infer1{\mbtin{E} \mbpeqca \mbtin{F}}
        \end{prooftree}
    \end{center}
  \end{minipage}

  \vspace{.5cm}

  \textbf{Type combinations} (\rocqlink{https://edgardschi.github.io/mailbox-types-rocq/MailboxTypes.Types.html\#TypeCombination}) \hfill \fbox{$\mbtcomb{T}{U}{T'}$}
  \begin{center}
    \begin{prooftree}
        \infer0{\mbtcomb{C}{C}{C}}
      \end{prooftree}
      \qquad
      \begin{prooftree}
        \infer0{\mbtcomb{\mbtout{E}}{\mbtout{F}}{\mbtout{(E \mbcomp F)}}}
      \end{prooftree}
      \qquad
      \begin{prooftree}
        \infer0{\mbtcomb{\mbtout{E}}{\mbtin{(E \mbcomp F)}}{\mbtin{F}}}
      \end{prooftree}
    \end{center}
    \begin{center}
      \begin{prooftree}
        \infer0{\mbtcomb{\mbtin{(E \mbcomp F)}}{\mbtout{E}}{\mbtin{F}}}
      \end{prooftree}
      \qquad
      \begin{prooftree}
        \hypo{\mbtcomb{T_1'}{T_2'}{T_3'}}
        \hypo{T_1 \mbpeqca T_1'}
        \hypo{T_2 \mbpeqca T_2'}
        \hypo{T_3 \mbpeqca T_3'}
        \infer4{\mbtcomb{T_1}{T_2}{T_3}}
      \end{prooftree}
    \end{center}
}

    \caption{The definitions of pattern and type permutations, as well as type combinations.}
    \label{fig:permutations-definition}
\end{figure}

% Before formally defining type combinations, we require some preliminary definitions.
In the original definitions of type combinations \cite{deliguoroMailboxTypesUnordered2018,fowlerSpecialDeliveryProgramming2023}, types are identified up to commutativity and associativity.
In practice, this means that types such as $\mbtout{(\mbmsg{m} \mbcomp \mbmsg{n})}$ and $\mbtout{(\mbmsg{n} \mbcomp \mbmsg{m})}$ are considered to be equivalent.
While this is true from a semantic point of view, syntactically these are different types
and assuming equality up to a certain property creates problems when mechanizing such definitions.
We require a way to express equivalence up to commutativity and associativity in
the context of type combinations.  To this end, we define a relation between
mailbox types, holding if they are permutations of each other.

The relation $E \mbpeqca F$ (Figure \ref{fig:permutations-definition}) is expressing that two patterns can be rewritten into each other using commutativity and associativity of the $\mbcomp$ operator.
% These rules are inspired by Rocq's definition of permutations of lists \cite{teamRocqStandardLibrary}, extended to be an equivalence relation.
Note how these rules do not include the $\mbchoice$ operator.
This is because type combination only considers type composition, and it is not required to include pattern choice in the definition of $\sim$ to achieve the desired properties.
Note the difference between the relations $E \sim F$ and $E \mbpeq F$.
The former is expressing that $E$ and $F$ can be rewritten into each other by swapping around patterns that appear between the composition operator.
The latter relates two semantically equal patterns.
For example, consider pattern $\mbmsg{m} \mbcomp \mbmsg{n}$.
Both $\mbmsg{m} \mbcomp \mbone \sim \mbone \mbcomp \mbmsg{m}$ and $\mbmsg{m} \mbcomp \mbone \mbpeq \mbone \mbcomp \mbmsg{m}$ hold,
while $\mbmsg{m} \mbcomp \mbone \mbpeq \mbmsg{m}$, but $\mbmsg{m} \mbcomp \mbone \not\sim \mbmsg{m}$.
% With pattern equality we are additionally able to express $\mbmsg{m} \mbcomp \mbone \mbpeq \mbmsg{m}$, but $\mbmsg{m} \mbcomp \mbone \not\sim \mbmsg{m}$.

% We now are able to extend $\mbpeqca$ to actual types, not just patterns.

% \begin{linkdef}{https://edgardschi.github.io/mailbox-types-rocq/MailboxTypes.Types.html\#TypeEq}
%   Two types are considered \emph{permutations} (written $T \mbpeqca U$) if it derivable using the following rules:
%
%   \begin{center}
%         \begin{prooftree}
%             \infer0{C \mbpeqca C}
%         \end{prooftree}
%         \qquad
%         \begin{prooftree}
%           \hypo{E \mbpeqca F}
%           \infer1{\mbtout{E} \mbpeqca \mbtout{F}}
%         \end{prooftree}
%         \qquad
%         \begin{prooftree}
%           \hypo{E \mbpeqca F}
%           \infer1{\mbtin{E} \mbpeqca \mbtin{F}}
%         \end{prooftree}
%     \end{center}
% \end{linkdef}

The relation $\sim$ is extended to mailbox types by relating the underlying patterns.
Base types are always considered to be permutations of each other.

% As was the case with some previous definitions, type combination is originally defined as a partial operator on types \cite{deliguoroMailboxTypesUnordered2018,fowlerSpecialDeliveryProgramming2023}.
% Again, instead of defining a function, we use a relation to capture the process of type combination.

% \begin{figure}[t]
%   % \textbf{Type combinations} \hfill \fbox{$\mbtcomb{T}{U}{T'}$}
%   \begin{center}
%     \begin{prooftree}
%         \infer0{\mbtcomb{C}{C}{C}}
%       \end{prooftree}
%       \qquad
%       \begin{prooftree}
%         \infer0{\mbtcomb{\mbtout{E}}{\mbtout{F}}{\mbtout{(E \mbcomp F)}}}
%       \end{prooftree}
%       \qquad
%       \begin{prooftree}
%         \infer0{\mbtcomb{\mbtout{E}}{\mbtin{(E \mbcomp F)}}{\mbtin{F}}}
%       \end{prooftree}
%     \end{center}
%     \begin{center}
%       \begin{prooftree}
%         \infer0{\mbtcomb{\mbtin{(E \mbcomp F)}}{\mbtout{E}}{\mbtin{F}}}
%       \end{prooftree}
%       \qquad
%       \begin{prooftree}
%         \hypo{\mbtcomb{T_1'}{T_2'}{T_3'}}
%         \hypo{T_1 \mbpeqca T_1'}
%         \hypo{T_2 \mbpeqca T_2'}
%         \hypo{T_3 \mbpeqca T_3'}
%         \infer4{\mbtcomb{T_1}{T_2}{T_3}}
%       \end{prooftree}
%     \end{center}
%     % TODO: Move this into the Pat chapter
%     % \begin{minipage}[t]{0.53\textwidth}
%     %   \begin{tabular*}{\linewidth}{@{}l@{\extracolsep{\fill}}r@{}}
%     %     \textbf{Usage-annotated} & \fbox{$\mbtucomb{A}{B}{A'}$}\\
%     %     \textbf{type combinations} &
%     %   \end{tabular*}
%     %   \begin{center}
%     %     \begin{prooftree}
%     %       \infer0{\mbtucomb{C}{C}{C}}
%     %     \end{prooftree}
%     %     \quad
%     %     \begin{prooftree}
%     %       \hypo{\mbtucomb{\eta_1}{\eta_2}{\eta}}
%     %       \hypo{\mbtcomb{J_1}{J_2}{K}}
%     %       \infer2{\mbtucomb{J_1^{\eta_1}}{J_2^{\eta_2}}{K^{\eta}}}
%     %     \end{prooftree}
%     %   \end{center}
%     % \end{minipage}
%     % \hfill
%     % \begin{minipage}[t]{0.46\textwidth}\vspace{-\baselineskip}
%     %   \textbf{Usage combinations} \hfill \fbox{$\mbtucomb{\eta_1}{\eta_2}{\eta_3}$}
%     %   \begin{align*}
%     %     \mbtucomb{\usecondclass}{\usecondclass}{\usecondclass}\\
%     %     \mbtucomb{\usecondclass}{\ureturnable}{\ureturnable}
%     %   \end{align*}
%     % \end{minipage}
%
%   \caption{The definition of type combination}
%   \label{fig:type-combination-definition}
% \end{figure}

% \begin{linkdef}{https://edgardschi.github.io/mailbox-types-rocq/MailboxTypes.Types.html\#TypeCombination}[Type Combination]
%   \emph{Type combination} is a ternary relation (written \mbtcomb{T}{U}{T'}), where for types $T$ and $U$, $T'$ is the type combination if it is derivable using the following rules:
%   % For types $T$ and $U$, $T'$ is the \emph{type combination} (written \mbtcomb{T}{U}{T'}) if it is derivable using the following rules:
%
%   \begin{center}
%     \begin{prooftree}
%         \infer0{\mbtcomb{C}{C}{C}}
%       \end{prooftree}
%       \qquad
%       \begin{prooftree}
%         \infer0{\mbtcomb{\mbtout{E}}{\mbtout{F}}{\mbtout{(E \mbcomp F)}}}
%       \end{prooftree}
%       \qquad
%       \begin{prooftree}
%         \infer0{\mbtcomb{\mbtout{E}}{\mbtin{(E \mbcomp F)}}{\mbtin{F}}}
%       \end{prooftree}
%     \end{center}
%     \begin{center}
%       \begin{prooftree}
%         \infer0{\mbtcomb{\mbtin{(E \mbcomp F)}}{\mbtout{E}}{\mbtin{F}}}
%       \end{prooftree}
%       \qquad
%       \begin{prooftree}
%         \hypo{\mbtcomb{T_1'}{T_2'}{T_3'}}
%         \hypo{T_1 \mbpeqca T_1'}
%         \hypo{T_2 \mbpeqca T_2'}
%         \hypo{T_3 \mbpeqca T_3'}
%         \infer4{\mbtcomb{T_1}{T_2}{T_3}}
%       \end{prooftree}
%     \end{center}
% \end{linkdef}


Type combinations model the concurrent interaction between processes acting on the same mailbox.
For example, assume processes $P_1$ and $P_2$ that both interact with mailbox $m$.
Process $P_1$ stores some message $\msg{m}$, while $P_2$ stores some message
$\msg{n}$.
Then, mailbox $m$ can be described by the type combination
$\mbtcomb{\mbtout{\mbmsg{m}}}{\mbtout{\mbmsg{n}}}{\mbtout{(\mbmsg{m} \mbcomp \mbmsg{n}})}$.
If instead, we have type $\mbtin{(\mbmsg{m} \mbcomp \mbmsg{n})}$, then we get the type combination
$\mbtcomb{\mbtout{\mbmsg{m}}}{\mbtin{(\mbmsg{m} \mbcomp \mbmsg{n}})}{\mbtin{\mbmsg{n}}}$.
Figure \ref{fig:type-combination} shows a example with three mailbox types.
There are different ways of arranging the combination, but in the end, the result always balances out to an empty mailbox.

\begin{figure}[t]
  \hbox{\scriptsize
  \hspace{-5pt}
  \begin{forest}
    for tree={grow'=90, s sep=0.2mm},
    for leaves={tier=leaf},
    [$\mbtin{\mbone}$
      [$\mbtin{(\mbmsg{n} \mbcomp \mbone)}$
        [$\mbtin{(\mbmsg{m} \mbcomp \mbmsg{n} \mbcomp \mbone)}$]
        [$\mbtout{\mbmsg{m}}$]
      ]
      [$\mbtout{\mbmsg{n}}$]
    ]
  \end{forest}
  \vline\quad\hspace{-3mm}
  \begin{forest}
    for tree={grow'=90, s sep=.2mm},
    for leaves={tier=leaf},
    [$\mbtin{\mbone}$
      [$\mbtin{(\mbmsg{m} \mbcomp \mbone)}$
        [$\mbtin{(\mbmsg{m} \mbcomp \mbmsg{n} \mbcomp \mbone)}$]
        [$\mbtout{\mbmsg{n}}$]
      ]
      [$\mbtout{\mbmsg{m}}$]
    ]
  \end{forest}
  \vline\quad\hspace{-3mm}
  \begin{forest}
    for tree={grow'=90, s sep=.2mm},
    for leaves={tier=leaf},
    [$\mbtin{\mbone}$
      [$\mbtin{(\mbmsg{m} \mbcomp \mbmsg{n} \mbcomp \mbone)}$]
      [$\mbtout{(\mbmsg{m} \mbcomp \mbmsg{n})}$
        [$\mbtout{\mbmsg{m}}$]
        [$\mbtout{\mbmsg{n}}$]
      ]
    ]
  \end{forest}
  }
  \caption{
    Example of possible type combinations for three mailbox types.
    Each top node represents the type of the same mailbox in three different processes that run in parallel.
    Edges indicate the type combination relation.
  }
  \label{fig:type-combination}
\end{figure}

Figure \ref{fig:permutations-definition} shows the definition of the \emph{type combination relation}.
Note how there is no rule for type combination of two types with input capabilities.
This would be allowing parallel reads of the same mailbox, which is not safe \cite{deliguoroMailboxTypesUnordered2018}, and thus not permitted.
The last rule in the definition of type combinations allows us to substitute types for their permutations.
Note how the previous definition of pattern equality $E \mbpeq F$ is specifically \emph{not} used in the definition of type combination.
%
The combination operation is a \emph{syntactic} operation, i.e.\ we can only
combine types that are \emph{syntactically equal} to the forms specified by the
type combination operation (up to commutativity and associativity).
To combine other mailbox types, for example
$\mbtout{\mbmsg{m}}$ and $\mbtin{(\mbmsg{m} \mbchoice \mbmsg{m})}$ we must firstly use
\emph{subtyping} to rewrite the latter mailbox type to $\mbtin{(\mbmsg{m} \mbcomp \mbone)}$
before using the combination operator. We will discuss subtyping further in
Section~\ref{sec:pat-programming-language}.


% SJF: think this has been said already
%%The need for a notion of equivalence up to certain properties is due to type combination being commutative and associative.

% TODO: Move this into the next Section
% Figure \ref{fig:type-combination-definition} shows the extension of type combination to \emph{usage-annotated type combination}.
% For usage annotations, we have to obey the rule that a $\usecondclass$ variable has to occur before a $\ureturnable$ variable.
% Otherwise, we would violate the property that any $\ureturnable$ variable appears as the final one in a process.
% Additionally, combining two $\ureturnable$ is not allowed, otherwise there could be multiple instances of a returnable mailbox name per thread.

\paragraph{Mechanization.}
Our mechanization of mailbox types differs from the original formalizations in that we treat partial operations as relations,
and include the notion of pattern and type permutations.
This addition allows us to prove desired properties, such as commutativity of type combinations.
Naturally, proofs about type combinations become more complex since we have to cover additional cases.
